\subsection{秩与有效性分析}
\label{sec: rank}


\begin{figure}[t]
    \centering
    \includegraphics[scale=0.36]{figure/comparision_with_rnn.png}
    \caption{PTB 数据集上测试集 PPL 随秩/隐单元数的变化。此处 RNNs 指 Vanilla-RNNs，其嵌入维度与 TTLM-Large 和 TTLM-Tiny 相同。RNNs-100、RNNs-200 与 RNNs-300 分别表示嵌入维度固定为 100、200 和 300 的 Vanilla-RNNs。}
    \label{fig:rank}
\end{figure}


TT 格式（TT format）的秩已被用于解释 RNNs 的表达能力或长期记忆容量 \citep{khrulkov2018expressive, levine2018benefits}。TT 分解的\textit{秩}已被证明等于 RNNs 隐状态的维度 \citep{khrulkov2018expressive}，这反映了 RNNs 的容量。秩越高，容量越大，反之亦然。然而，秩与语言建模有效性之间的关系尚未在实践中得到展示。我们将针对秩评估 TTLM-Large 与 TTLM-Tiny 的有效性。

\subsubsection{有效性}
表 \ref{tab:evaluation} 给出了我们的模型与各基线模型在 WikiText-2 和 PTB 数据集上的测试集 PPL 结果。如图所示，与 Vanilla-RNNs 相比，TTLM-Large 分别将 PPL 降低了 14.3 和 16.0，TTLM-Tiny 分别将 PPL 降低了 1.7 和 8.5。因此，当隐单元数或秩设为 20、嵌入维度为 400 时，TTLM-Large 与 TTLM-Tiny 的表现均优于所有基线模型。
\begin{figure}[t]
    \centering
    \includegraphics[scale=0.24]{figure/ICML/rank_val_ppl_Large_Tiny_ICML.png}
    \caption{TTLM-Large 与 TTLM-Tiny 在 PTB 数据集上随秩增大时的验证集 PPL。上：TTLM-Large，下：TTLM-Tiny。}
    \label{fig:rank_analysis}
\end{figure}

为进一步评估我们模型的有效性，我们在秩逐渐增大的条件下对 TTLM-Large、TTLM-Tiny 与 Vanilla-RNNs 进行了比较，如图 \ref{fig:rank} 所示。值得注意的是，我们还纳入了 RNNs-100、RNNs-200 与 RNNs-300，以控制较大嵌入维度可能带来的影响。如图所示，即使隐单元数达到 40，RNNs 的测试集 PPL 仍在持续下降，而 TTLM-Large 与 TTLM-Tiny 并非如此。因此，我们期望我们的模型在\textit{小规模}的隐单元（即数量在 5 到 40 之间）下优于 Vanilla-RNNs，但在更大规模下则不然。

\subsubsection{过拟合}
\begin{figure*}[t]
    \centering
    % \vspace{-15pt}
    \includegraphics[width= 1 \textwidth]{figure/four_comparison.png}
    \caption{非线性对 PTB 数据集上各 TTLM 变体的影响。后缀 \texttt{-tanh} 表示使用 $tanh$ 激活函数的模型，以虚线标示。Second-linear 指不带激活函数的二阶 RNNs。设置：所有模型均有 17 个隐单元/秩。}
    \label{fig:activation_comparision}
\end{figure*}
图 \ref{fig:rank_analysis} 展示了随秩数增大，TTLM-Large 与 TTLM-Tiny 在验证集上的表现。如图所示，当我们逐渐增大 TTLM-Large 的秩时，其验证集 PPL 在较早的训练轮次便开始上升。相比之下，TTLM-Tiny 的验证集 PPL 随秩数的增大而稳定下降。这一比较表明，TTLM-Large 比 TTLM-Tiny 更容易过拟合。图 \ref{fig:rank} 中的结果进一步支持了这一发现。TTLM-Tiny 的测试集 PPL 随秩的增大持续改善，而 TTLM-Large 的性能在秩数达到 25 时开始下降。

当我们关注两个模型之间的差异时，TTLM-Large 拥有一个额外的参数张量 $\tW^{eh}$。因此我们认为，TT 核（TT cores）的参数化越简单，模型越容易避免过拟合。这一发现与 \cite{wu_multiplicative_2016} 中 MI-RNNs 与二阶 RNNs 的比较结果一致。就实际情形而言，我们需要意识到，与 TTLM-Large 相比，TTLM-Tiny 拟合训练数据的能力较低，因而过拟合的风险也较低。


\subsection{非线性分析}
\label{sec: activation}
先前的研究曾尝试将 TT 分解作为理论平台来研究 RNNs \citep{khrulkov2018expressive, levine2018benefits}。然而，TT 分解与现有神经网络（如 Vanilla-RNNs）之间的一个关键差异，在于网络模型内部的非线性激活函数。张量分解中非线性的缺失使人质疑其理论分析能否迁移到基于 RNNs 的模型上。为理解 $tanh$ 激活函数对 TT 变体的作用是否随 TT 核的不同而变化，我们给出了如图 \ref{fig:activation_comparision} 所示的实证结果。


就收敛速度而言，$\tanh$ 加速了 TTLM-Large-tanh、TTLM-Tiny-tanh 与 MI-RNNs，而对二阶 RNNs 几乎没有影响。就最低验证集 perplexity 的量级而言，$\tanh$ 损害了 TTLM-Large 与 TTLM-Tiny 的性能，但对乘性积分（multiplicative integration）以及二阶 RNNs 中三阶张量 $\tT$ 的影响甚微。

因此，非线性激活函数对 TTLM 变体的影响取决于 TT 核的设置，无论是对验证集 PPL 的收敛还是对最低验证集 PPL 的量级而言都是如此。从实验的角度看，我们认为非线性函数对某一 TT 变体的作用不能简单地迁移或类推到另一个 TT 变体上。这也提示我们，对先前研究所宣称的张量分解与现有神经网络模型在实现层面的类比应保持警惕 \citep{khrulkov2018expressive, levine2018benefits}。非线性激活函数可能是影响这种类比的一个因素。

